\chapter{Discussion}
\label{sec:discussion}

The results support a cautious but useful interpretation of prompt
decomposition for telematics graph recovery. The best observed configuration
is a decomposition-style prompt: \texttt{deepseek-r1:8b} with
\strategy{Strategy_26_Skeleton_of_Thought}. The broader evidence is not a
simple story that more structure always helps. Some structured strategies
produce empty outputs or many false positives, and the consensus graph
preserves recurring mistakes as well as recurring correct edges.

\section{Answering the Research Questions}

\subsection{RQ1: Model Ability to Recover Directed Relationships}

Both models recover some directed telemetry relationships, but neither
recovers the six-edge ground truth reliably. The strongest cell recovers four
of six ground-truth edges and predicts five additional false edges. Llama's
best cell recovers three of six ground-truth edges and also predicts six false
edges. These values show that local open-source LLMs can propose useful graph
fragments, especially for semantically salient energy and motion relations,
but they do not produce a sufficiently reliable graph without evaluation.

The most consistently recovered true edge is the energy-to-range relationship
from \texttt{electricenergyremaining} to \texttt{electricremaining}, which
appears in seven of the ten predicted edge sets. This is unsurprising because
the feature names strongly suggest a direct relation. Edges that require more
careful direction or directness judgment are less stable.

\subsection{RQ2: Prompt Decomposition Strategies}

Prompt decomposition has mixed results. Skeleton-of-thought is the strongest
strategy by mean F1 score across the two models and produces the best single
cell. Its subsystem-oriented skeleton appears to help DeepSeek connect
energy-state, driving-behaviour, and cross-group relationships. This is the
clearest positive evidence for decomposition in the thesis.

Least-to-most prompting does not show the same benefit. DeepSeek's
least-to-most sidecar contains an empty raw response and no parsed edges,
while Llama's least-to-most output contains many false positives. This
suggests that decomposition can fail in two different ways: it can make the
model fail to produce a parseable answer, or it can encourage a broad list of
plausible but unsupported links. Decomposition is therefore not a guarantee of
better graph recovery; it is a design variable that must be validated per
model and task.

\subsection{RQ3: Instruction Structuring}

Policy-first structured prompting performs weakly in the validated results.
DeepSeek produces no parseable edges under this strategy, and Llama recovers
only one true edge while producing five false positives. The policy-first
prompt is methodologically reasonable because it emphasizes directness,
interventionability, and output format, but the result shows that explicit
rules alone are not sufficient. A model can still omit the graph, reverse
directions, or include plausible non-ground-truth links.

This finding is useful because it counters a common assumption in LLM system
design: stating a policy before the task does not automatically improve
output quality. For graph recovery, constraints must be paired with a prompt
shape that helps the model cover the right feature groups while preserving the
final edge-list format.

\subsection{RQ4: Strongest Model-Strategy Combination}

The strongest observed combination is \texttt{deepseek-r1:8b} with
skeleton-of-thought prompting. It achieves F1 score 0.5333, recall 0.6667,
precision 0.4444, SHD 7, and Pearson correlation 0.4303. This combination
recovers four of the six ground-truth edges. It is the only row with F1 above
0.4000.

The result should not be generalized beyond the experiment. It is one output
from one model under one prompt strategy on one Tesla EV graph. It does,
however, show that a subsystem-based skeleton can improve graph recovery for
this specific DeepSeek configuration.

\subsection{RQ5: Recurring Correct and Incorrect Relationships}

The consensus graph shows that repeated agreement across model-strategy cells
does not equal correctness. The consensus edge
\texttt{electricenergyremaining -> electricremaining} is correct and has
support 0.7. The edge \texttt{acceleration\_kmh2 -> speed\_kmh} is also
correct and has support 0.5. However, the consensus graph also includes
\texttt{speed\_kmh -> acceleration\_kmh2}, \texttt{mileage -> distance\_m},
and \texttt{distance\_m -> mileage}, all of which are false under the
ground-truth DAG.

This pattern shows that the models often identify related feature pairs but
struggle with direction and directness. In a fleet-monitoring setting, this is
a serious failure mode: a graph can appear plausible while encoding the wrong
direction of influence.

\section{Why Skeleton-of-Thought Helped}

The skeleton-of-thought prompt fits the structure of the Tesla EV task. It
asks the model to consider energy-state edges, driving-behaviour edges, and
cross-group edges before applying an integrity gate. The ground-truth DAG also
has this structure: one relationship is within energy state, two involve
driving behaviour and distance/usage, and two connect driving behaviour to
remaining energy. This alignment likely explains why DeepSeek with
skeleton-of-thought recovers the most ground-truth edges.

The same prompt is less successful for Llama. Llama skeleton output contains
two true positives and seven false positives. The model covers many related
feature pairs, but it overconnects mileage, distance, and energy variables.
The prompt therefore appears to improve coverage more than precision for
Llama. This difference reinforces a central point of the thesis: prompt
structure interacts with model behavior.

\section{Why Structured Prompts Failed}

Structured prompts fail when they optimize the wrong behavior. A policy-first
prompt can make the rules explicit but still fail to generate a useful graph.
A least-to-most prompt can encourage broad causal brainstorming rather than
strict final edge selection. A skeleton prompt can improve coverage but also
encourage extra cross-group links. These failures are visible because the
evaluation parses outputs into matrices and counts false positives, false
negatives, and structural disagreement.

The DeepSeek empty-output cases are especially important. The sidecar JSON
files for DeepSeek zero-shot, policy-first, and least-to-most contain empty
raw responses and zero parsed edges, while reporting no errors. This is a
pipeline-level fragility: a local model call may complete but still produce no
usable edge list. Such cells receive low recall and F1, but they also receive
FDR 0.0000 because no false edges were predicted. This is why FDR must be read
with recall and F1.

\section{DeepSeek Versus Llama3.1}

DeepSeek is more uneven in the validated artifact set. It produces the
best graph under skeleton-of-thought and a competitive graph under the
no-instruction control, but it also produces three empty outputs. Llama is
more consistent in producing edge lists, but its outputs are dense and
false-positive heavy. This is reflected in the model averages: Llama has
higher mean F1, while DeepSeek has lower mean SHD and much lower mean FDR.

This does not support a universal model ranking. It supports a deployment
lesson: the model that produces the best single graph is not necessarily the
most dependable across prompt conditions. In a real fleet-monitoring workflow,
a system designer would need both quality controls for empty outputs and
filters for false-positive edges.

\section{Consensus Graph Interpretation}

The consensus graph is useful as a diagnostic artifact, not as an empirical
replacement for the ground truth. It shows which edges are attractive to the
models across prompt strategies. Some of those edges are correct; others are
recurring misconceptions. The false consensus edge
\texttt{speed\_kmh -> acceleration\_kmh2} is particularly instructive because
it is the reverse of a ground-truth edge. The models appear to treat the two
signals as tightly related but are inconsistent about causal direction.

The mileage-distance edges show a similar issue. The models frequently link
\texttt{mileage} and \texttt{distance\_m} in one or both directions, but the
ground truth encodes speed as the source for both mileage and distance. This
suggests that models may compress accumulated and short-window distance
signals into a generic distance relation rather than respecting the specific
graph definition.

\section{Implications for Fleet Monitoring}

The results support using LLMs as graph proposal tools only when their outputs
are evaluated. A proposed graph can help analysts inspect possible telemetry
relationships, but it should not directly drive maintenance decisions,
diagnostics, or control logic without validation. The value of the pipeline is
that it turns a free-form LLM output into a measurable artifact.

For fleet-monitoring applications, the practical workflow should therefore be:
define the feature set, provide a compact data fingerprint, constrain the
output format, parse the proposed graph, compare it against known or
expert-reviewed structure, and inspect false positives. Prompt decomposition
can be useful in this workflow, but it is not a substitute for graph
validation.

\section{Relation to Earlier Repository Infrastructure}

Earlier repository work remains relevant as engineering context. It
established data-minimization habits, local inference through Ollama, and
artifact traceability. In this thesis, those elements support the empirical
workflow rather than define the central claim. The core result is about graph
recovery from telematics data, with governance discussed where it affects
deployment practice.
